% !PRINTED 391
\weberpage{391}

\S.\,104.

\gesperrt{Realität der Doppeltangenten.}

Wir wollen noch die Frage erörtern, wie sich bei einer
reellen Curve vierter Ordnung ohne singulären Punkt die Doppeltangenten
in Bezug auf ihre Realität verhalten können. Wir
nehmen also jetzt die Co\"efficienten der Gleichung der Curve
vierter Ordnung reell an, d.\,h. wir setzen einen reellen Rationalitätsbereich
voraus.

Wenn dann eine Doppeltangente $\xi$ imaginär ist, so muss
auch die conjugirte Gerade $\xi'$ Doppeltangente sein.

Wenn wir in einer rationalen Gleichung zwischen den Wurzeln
der Doppeltangentengleichung jede imaginäre Wurzel durch
die conjugirte ersetzen, so entsteht wieder eine richtige Gleichung.
Es folgt daraus, dass zu syzygetischen oder azygetischen
Systemen vom % [sic] \textquotedbl{}vom\textquotedbl{} printed here where \textquotedbl{}von\textquotedbl{} would be expected
Doppeltangenten conjugirte Systeme desselben
Charakters gehören, die man erhält, wenn man überall $i$ durch
$-i$ ersetzt.

Ein \gesperrt{Steiner}'scher Complex $C$ geht daher durch Uebergang
von $i$ zu $-i$ in einen \gesperrt{Steiner}'schen Complex $C'$ über und wir
unterscheiden zwei Fälle:

Ist $C$ mit $C'$ identisch, so nennen wir $C = C'$ einen \gesperrt{reellen
Complex}.

Ist aber $C$ von $C'$ verschieden, so bilden sie ein \gesperrt{conjugirtes
Complexpaar}.

Die Paare eines reellen Complexes bestehen entweder aus
zwei reellen Doppeltangenten oder es sind conjugirte Paare $\xi, \xi'$,
oder sie enthalten zwei nicht conjugirte imaginäre Doppeltangenten
$\xi, \eta$. Im letzteren Falle muss dann im Complex auch
das aus den conjugirten Doppeltangenten $\xi', \eta'$ gebildete Paar
vorkommen. Zwei solche Paare $(\xi, \eta)$, $(\xi', \eta')$ wollen wir ein
\gesperrt{conjugirtes Doppelpaar} nennen.

Niemals kommt in einem reellen Complexe eine reelle
Wurzel $x$ mit einer imaginären $\xi$ gepaart vor, weil sonst neben
dem Paare $(x, \xi)$ auch das conjugirte Paar $(x, \xi')$ in demselben
Complexe vorkommen müsste, was unmöglich ist, da diese beiden
Paare ein gemeinschaftliches Element $x$ enthalten würden.

Ein \gesperrt{Aronhold}'sches Siebener-System $S$ geht durch Vertauschung
aller seiner Elemente mit den conjugirten in ein eben-
%CONT>>
